\chapter{Background and Literature Review}\label{ch:literature}

\section{Chess Engines and Computer Chess}

Classical chess engines combine alpha-beta search with an evaluation function. The landmark
1997 match in which IBM's Deep Blue defeated Garry Kasparov showed that hand-crafted evaluation
combined with massive search could beat the world champion. Modern engines retain that search
architecture but have replaced the hand-written evaluation. Stockfish~\cite{stockfish}, the
strongest open-source engine, evaluates positions with an \emph{efficiently updatable neural
network} (NNUE), originally developed for computer shogi by Nasu~\cite{nasu2018nnue}. Because the
network's first layer can be updated incrementally as pieces move, a neural evaluation becomes
cheap enough to run inside a search visiting millions of positions per second. That property
matters for this dissertation, which analyses about 18,000 games.

A second line of work replaces search heuristics with learning altogether. AlphaZero learns a
policy and a value function by self-play and combines them with Monte Carlo tree search, reaching
superhuman strength from the rules alone~\cite{silver2017mastering,silver2018general}. Zahavy et
al.~\cite{zahavy2023diversifying} show that a single AlphaZero policy is near-optimal but narrow. It
fails on out-of-distribution positions, such as Penrose-style puzzles that are easy for humans.
Their AZdb system trains a population of latent-conditioned policies with diversity bonuses; it
solves more hard puzzles and plays in more varied styles.

These systems are designed to \emph{play} well. Maia~\cite{mcilroyyoung2020aligning} takes the
opposite view: it trains networks on human games to predict the move a human \emph{of a given
rating} would play, rather than the best move. For a training system this distinction is central.
The engine tells us what is correct; a human-behaviour model would tell us what a particular player
is likely to find. This dissertation uses Stockfish as an oracle of correctness. Maia-based
difficulty estimates are discussed as future work (Section~\ref{sec:future}).

\section{Chess Puzzles and Tactics Training Platforms}

The \textbf{Lichess puzzle generator}~\cite{lichesspuzzler} scans games from the Lichess database
with Stockfish, finds positions where one side has a decisive and unique continuation, and tags each
puzzle with themes such as \emph{fork}, \emph{pin} or \emph{backRankMate}. The resulting open
database~\cite{lichessdb} contains about six million puzzles, each rated with Glicko-2 from the
attempts of Lichess users. It is the main puzzle source for this project, and its themes serve as
the reference labels in Section~\ref{sec:labeller-results}.

\textbf{Chessity}~\cite{chessity} offers hand-crafted tactics and adapts difficulty to performance,
keeping players challenged without discouraging them. \textbf{Aimchess}~\cite{aimchess} and
\textbf{Noctie}~\cite{noctie} analyse a user's own games and build training material from their most
instructive mistakes. Nie et al.~\cite{nie2026discovering} study the pedagogical value of puzzles
directly. Using 1.5 billion puzzle-solving interactions, they apply offline reinforcement learning
and off-policy evaluation to recommend puzzles, and report the largest benefit for beginners whose
progress had stalled. Table~\ref{tab:tools} positions this project against these systems.

\begin{table}[H]
\centering
\small
\caption{Related puzzle and training systems.}\label{tab:tools}
\begin{tabularx}{\linewidth}{p{0.24\linewidth}Yp{0.13\linewidth}p{0.13\linewidth}}
\toprule
\textbf{System} & \textbf{Source of puzzles} & \textbf{Adaptive} & \textbf{Own games} \\
\midrule
Lichess puzzler~\cite{lichesspuzzler} & Engine-extracted from community games & No & Indirectly \\
Feng et al.~\cite{feng2025generating} & Generative model and RL fine-tuning & No & No \\
Chessity~\cite{chessity} & Hand-crafted, fixed pool & Difficulty & No \\
Aimchess / Noctie~\cite{aimchess,noctie} & Mined from the user's games & Mistake-based & Yes \\
Nie et al.~\cite{nie2026discovering} & Lichess pool, offline-RL selection & Yes (learned) & No \\
\textbf{This work} & Lichess pool and engine-verified mining from own games & Category \emph{and}
difficulty, online Bayesian & Yes \\
\bottomrule
\end{tabularx}
\end{table}

\section{Player Modelling and Decision-Making Style}

Can a player's ``style'' be measured? McIlroy-Young et al.~\cite{mcilroyyoung2021detecting} introduce
\emph{behavioural stylometry} for chess. A transformer over move sequences embeds a set of games
into a style space, and a player can be identified among thousands of candidates from 100 labelled
games with 98\% accuracy. The embeddings transfer from club players to grandmasters. The authors
point out the privacy risk: a playing style can act as a behavioural fingerprint.

Jayasekara~\cite{jayasekara2018classification} asks whether classical style labels (aggressive,
positional, defensive) can be recovered without supervision. Using engine-derived features from the
opening and middlegame, with PCA and Gaussian mixture models, some styles separate clearly and others
(for example ``solid'' and ``tactical'') merge.

Together these works treat style as latent structure in high-dimensional data rather than as a
metaphor. This dissertation models the player with more modest, directly interpretable quantities:
overall ability, per-category tactical offsets, and a rating-group prediction from engine-derived
features. One of its findings is that per-category tactical differences between players are much
harder to detect than overall strength (Section~\ref{sec:cohort-results}), which is itself relevant
to the style-modelling literature.

\section{Generative Models and Reinforcement Learning}

\paragraph{Generative puzzle creation.} Feng et al.~\cite{feng2025generating} pre-train a discrete
generative model on 4.4 million Lichess puzzles and fine-tune it with reinforcement learning. The
reward is built from engine search statistics that measure uniqueness, counter-intuitiveness (the
key move is hard for a shallow search to find), diversity and realism, and explicit filters prevent
reward hacking. RL fine-tuning raises the rate of counter-intuitive puzzles tenfold, and human
experts rated the output as more creative than composed book puzzles. This is the state of the art
the portfolio aimed at. It is also the part of the plan this dissertation did \emph{not} implement
(Section~\ref{sec:generation}).

\paragraph{Bandits as reinforcement learning.} Choosing which category to practise next, from
binary feedback, is a multi-armed bandit problem, the simplest setting of the exploration and
exploitation trade-off in reinforcement learning~\cite{sutton2018reinforcement}. Thompson's rule of
1933~\cite{thompson1933likelihood} keeps a posterior for each arm, draws one sample from each, and
acts greedily on the samples. Exploration comes from posterior uncertainty rather than from a tuned
parameter. Upper-confidence-bound methods achieve logarithmic regret through deterministic
optimism~\cite{auer2002finite}. Chapelle and Li~\cite{chapelle2011empirical} found Thompson Sampling
competitive with or better than UCB in practice, and Russo et al.~\cite{russo2018tutorial} give the
modern analysis. When rewards change over time, discounted~\cite{raj2017taming} and
sliding-window~\cite{garivier2011upper} variants forget old evidence.

In education, Clement et al.~\cite{clement2015multi} used bandits to sequence exercises for learners,
rewarding measured \emph{learning progress} rather than raw success. That distinction reappears in
Chapter~\ref{ch:results}.

\section{Difficulty Modelling and Educational Outcomes}

Elo's rating system~\cite{elo1978rating} models a game's outcome as a logistic function of the
rating difference. Glicko~\cite{glickman1999parameter} adds a \emph{rating deviation} that expresses
uncertainty. Glicko-2~\cite{glickman2012example} adds a \emph{volatility} that expresses how
erratically true strength moves. Lichess rates its puzzles with Glicko-2, treating each attempt as a
game between the solver and the puzzle. Pelánek~\cite{pelanek2016applications} surveys the use of
Elo-type models in adaptive educational systems. They estimate learner ability and item difficulty
on one scale and update online, which is exactly the role they play in this project. The same
logistic form is the Rasch model of item response theory.

What makes a tactical problem hard for a human is not obvious. Hristova et
al.~\cite{hristova2014assessing} find only a weak relation between players' strength and their
ability to judge problem difficulty. Their eye-tracking results suggest that difficulty comes mainly
from the depth and reliability of calculation, not from recognising the motif. This caution is
relevant here, because the system's taxonomy is organised by motif.

On the educational side, knowledge tracing~\cite{corbett1995knowledge} maintains a per-skill
probability of mastery, updated by Bayes' rule from correct and incorrect answers. It is the direct
ancestor of the per-category posteriors used here. Vygotsky's zone of proximal
development~\cite{vygotsky1978mind} and Ericsson's account of deliberate practice~\cite{ericsson1993role}
supply the pedagogical rationale. Material should be hard enough to require effort but not so hard
that it only produces failure, and it should concentrate on current weaknesses.

\paragraph{Predicting puzzle difficulty from content.} Estimating how hard a puzzle is \emph{before}
anyone has attempted it became a shared task in 2024: the IEEE Big Data Cup asked entrants to predict
Lichess puzzle ratings from the board and the solution alone, over a database of more than four
million puzzles~\cite{zysko2024bigdatacup}. The published entries are neural: a network whose
architecture imitates stages of human problem solving~\cite{schutt2024estimating}, and a transformer
that predicts the Glicko-2 rating from a spatio-temporal encoding of position and moves~\cite{milosz2024transformers}.
Earlier work in the cognitive-science tradition instead derived difficulty from properties of the search
tree, such as how deep and how branching the winning line is~\cite{hristova2014assessing,bratko2021predicting}.
What none of this work does is feed the estimate back into a training system: difficulty is predicted as
an end in itself, and reported as a regression error. The present work reuses the same public labels but
requires the estimate to carry a usable uncertainty, because a difficulty with no error bar cannot enter
a Bayesian ability update in a principled way (Section~\ref{sec:puzzlenet}).

\section{Summary and Research Gap}

Existing systems cover complementary parts of the problem. Lichess shows that engine-based
extraction scales. Feng et al.\ show that generative RL can create strong puzzles, but not for a
particular learner. Chessity adapts difficulty over a fixed pool. Aimchess and Noctie mine the
learner's own games. Nie et al.\ learn puzzle selection offline from enormous interaction logs.

The portfolio identified the gap as learner-conditioned \emph{generative} puzzle creation. That gap
remains. This dissertation addresses a narrower but related one that is more open to rigorous
evaluation. It combines online, uncertainty-aware selection of \emph{both} tactical category
\emph{and} difficulty for an individual learner, starting from their real games. It then tests,
against real data, whether the assumption behind every mistake-based trainer holds: that a
player's tactical weaknesses can be diagnosed from their games.
